\section{Solutions for Chapter~\ref{ch:cft}}

\paragraph{Exercise~\ref{ex:cft-level1}.}
\phantomsection
\label{sol:ex:cft-level1}
The commutator $[L_1,L_{-1}]=2L_0$ gives, by induction,
\[
 L_1L_{-1}^n|h\rangle
 =n(2h+n-1)L_{-1}^{n-1}|h\rangle.
\]
Taking an inner product and using $L_{-1}^\dagger=L_1$
therefore multiplies the previous norm by $n(2h+n-1)$.
Starting with a unit primary gives $n!(2h)_n$.
For equal external weights, the primary three-point function has
$\langle\phi(\infty)\phi(1)\mathcal O_p(t)\rangle
=C(1-t)^{-p}$. Differentiating $n$ times at $t=0$ inserts
$L_{-1}^n\mathcal O_p$ and gives $C(p)_n$.
Each three-point descendant coefficient is therefore $(p)_n$
relative to the primary coupling.
Writing $x$ for the cross-ratio as in the chapter, the block in the
reduced correlator $x^{2h}G(x)$ is consequently
\[
 x^p\sum_{n\ge0}\frac{(p)_n^2}{n!(2p)_n}x^n
 =x^p\left(1+\frac p2x+
       \frac{p(p+1)^2}{4(2p+1)}x^2+\cdots\right).
\]
Multiplying by $x^{-2h}$ restores the block of the unreduced
correlator displayed in the chapter.
The identity block is treated separately at $p=0$.



\paragraph{Exercise~\ref{ex:cft-null}.}
\phantomsection
\label{sol:ex:cft-null}
Let $v=(L_{-2}+aL_{-1}^2)|h\rangle$. Acting with $L_1$
gives $[3+a(4h+2)]L_{-1}|h\rangle$, so at $h=1/16$
we need $a=-4/3$. The $L_2$ action is
$4h+c/2+6ah$. At $c=1/2,h=1/16$ it equals
$1/4+1/4-1/2=0$. Thus
\[
 \left(L_{-2}-\frac43L_{-1}^2\right)|1/16\rangle
\]
is the null vector. Higher positive modes annihilate it by
grading. Its vanishing is the input for the Ising differential
equation in the sewing chapter.



\paragraph{Exercise~\ref{ex:cft-compact}.}
\phantomsection
\label{sol:ex:cft-compact}
For $(m,n)=(1,1)$, $\bar h=0$ requires $R^{-1}-R/2=0$,
hence $R=\sqrt2$. Then $p_L=\sqrt2$ and $h=p_L^2/2=1$.
Under $(R,m,n)\mapsto(2/R,n,m)$, $p_L$ is unchanged
and $p_R$ changes sign. Both conformal weights are unchanged.
The holomorphic current at the fixed radius is the exponential
current of the $A_1$ lattice example.



\paragraph{Exercise~\ref{ex:cft-stress-combination}.}
\phantomsection
\label{sol:ex:cft-stress-combination}
The fourth-order pole of $(T_1+T_2)(z)
(c_2T_1-c_1T_2)(0)$ is
$(c_2c_1-c_1c_2)/2=0$, and the remaining singular poles
are $2A/z^2+\partial A/z$. Thus $A$ is primary for the
total stress. Orthogonality of the two stress states gives
\[
 \norm A^2=\frac12c_1c_2(c_1+c_2)>0.
\]
Its self-OPE has second-order coefficient
$2(c_2^2T_1+c_1^2T_2)$. Equality with $2A$ would require
$c_2^2=c_2$ and $c_1^2=-c_1$. The second condition is impossible
for $c_1>0$, so $A$ is not a stress tensor in this normalization.
A primary of weight two is therefore not automatically a
Virasoro stress.



\paragraph{Exercise~\ref{ex:cft-thermal}.}
\phantomsection
\label{sol:ex:cft-thermal}
Put $t=v\beta/L$. The vacuum contribution to $Z(it)$ at large
$t$ is $e^{\pi ct/6}$, because the cylinder vacuum energy
uses $c_L+c_R=2c$. Modular invariance replaces $t$ by $1/t$.
As $\beta\to0$,
\[
 \log Z\sim\frac{\pi c}{6t}=\frac{\pi cL}{6v\beta}.
\]
If the vacuum is not normalizable, the low-temperature trace is
controlled by the actual spectral threshold and its measure.
One cannot insert a unit-degeneracy vacuum term; this is where
effective central charge and possible power or logarithmic
prefactors enter.



\paragraph{Exercise~\ref{ex:cft-marginality}.}
\phantomsection
\label{sol:ex:cft-marginality}
One needs normalized OPE coefficients controlling quadratic beta
functions; higher-order integrated correlators, including local
contributions when insertion points coincide, controlling further
obstructions; and an analysis of which couplings describe inequivalent
theories while preserving the specified symmetry.
The coincident-point contributions are called \emph{contact terms}.
Couplings removable by a field redefinition or a symmetry
transformation are \emph{redundant directions} and must not be
counted as independent parameters. One must also check whether a
putative marginal field becomes a descendant when a current ceases
to be conserved; this is one mechanism of \emph{multiplet
recombination}. These are additional tasks, not data supplied by
the character count.
Vanishing beta functions must hold on a locus of couplings after
these identifications. A count of thirty-six dimension-two
primaries alone supplies none of those conclusions.

